% ---
% school: "college"
% level: "8eme"
% section: null
% chapter: "التناظر المركزي"
% source: "Axiomaths"
% title: "صحيح أم خاطئ ؟"
% country: "Tunisie"
% year: 2026
% professor: "Axiomaths"
% difficulty: "Facile"
% points: null
% tags: ["symetrie-centrale", "vrai-faux", "proprietes", "contre-exemple"]
% ---
% Exercice 3/15 de la série 8ème « التناظر المركزي » (statement + \axsol + solution). Compiler avec XeLaTeX + axiomath.sty (+ axiomaths-logo.png).
\documentclass[12pt,a4paper]{article}
\usepackage[level=8b,color=15803D]{axiomath}
\usetikzlibrary{arrows.meta}
\newcommand{\cm}{\,\mathrm{cm}}
\axsetup{title={التناظر المركزي}, corrige}

\begin{document}

\begin{corr}[title={صحيح أم خاطئ ؟}, stars=1, source={Axiomaths}]
في كل حالة، حدّد إن كانت العبارة صحيحة أم خاطئة ثم علّل إجابتك.
\begin{enumerate}
  \item مناظر قطعة $[AB]$ بالنسبة إلى نقطة $O$ هي قطعة $[A'B']$ حيث $A'B'=AB$.
  \item إذا كانت $O$ منتصف $[MN]$ فإن $N$ مناظرة $M$ بالنسبة إلى $O$.
  \item مناظر مستقيم بالنسبة إلى نقطة هو مستقيم عمودي عليه.
  \item $O$ منتصف $[AB]$ و $C$ مناظرة $A$ بالنسبة إلى $B$. إذن $O$ منتصف $[AC]$.
  \item مناظرة دائرة شعاعها $r$ بالنسبة إلى نقطة هي دائرة شعاعها $2r$.
  \item مناظرة زاوية قيسها $50^\circ$ بالنسبة إلى نقطة هي زاوية قيسها $50^\circ$.
\end{enumerate}
\axsol
\begin{enumerate}[itemsep=10pt]
  \item \textbf{صحيح.} التناظر المركزي يحافظ على الأطوال.
  \item \textbf{صحيح.} هذا هو تعريف المناظرة : $O$ منتصف $[MN]$ معناه $N=S_O(M)$.
  \item \textbf{خاطئ.} مناظر المستقيم مستقيم \emph{مواز} له (أو منطبق عليه).
  \item \textbf{خاطئ.} مثال مضاد : $AB=4\cm$ ؛ $O$ على بعد $2\cm$ من $A$ ، و $C$ على بعد $8\cm$ من $A$ ، فمنتصف $[AC]$ على بعد $4\cm$ من $A$ وليس $O$. في الحقيقة $B$ هي منتصف $[AC]$.
  \item \textbf{خاطئ.} الدائرة المناظرة لها نفس الشعاع $r$ (الأطوال محفوظة) ؛ مركزها هو مناظر المركز.
  \item \textbf{صحيح.} التناظر المركزي يحافظ على قيس الزوايا.
\end{enumerate}
\end{corr}

\end{document}
